%%%%%%%%%%%%%%%%%%%%%%%%%%%%%%%%%%%%%%%%%%%%%%%%%%%%%%%%%%%%%%%%%%%%%%%%%
%
% Apuntes de la asignatura Análisis Matemático II.
% Doble Grado de Informática y Matemáticas.
% Universidad de Granada.
% Curso 2016/17.
% 
% 
% Colaboradores:
% Javier Sáez (@fjsaezm)
% Daniel Pozo (@danipozodg)
% Pedro Bonilla (@pedrobn23)
% Guillermo Galindo (@guillegalor)
% Antonio Coín (@antcc)
% Sofía Almeida (@SofiaAlmeida)
% Miguel Lentisco (@AlfaOmegaX)
%
% Agradecimientos:
% Andrés Herrera (@andreshp) y Mario Román (@M42) por
% las plantillas base.
%
% Sitio original:
% https://github.com/libreim/apuntesDGIIM/
%
% Licencia:
% CC BY-NC-SA 4.0 (https://creativecommons.org/licenses/by-nc-sa/4.0/)
%
%%%%%%%%%%%%%%%%%%%%%%%%%%%%%%%%%%%%%%%%%%%%%%%%%%%%%%%%%%%%%%%%%%%%%%%%%


%------------------------------------------------------------------------------
%   ACKNOWLEDGMENTS
%------------------------------------------------------------------------------

%%%%%%%%%%%%%%%%%%%%%%%%%%%%%%%%%%%%%%%%%%%%%%%%%%%%%%%%%%%%%%%%%%%%%%%%
% Plantilla básica de Latex en Español.
%
% Autor: Andrés Herrera Poyatos (https://github.com/andreshp) 
%
% Es una plantilla básica para redactar documentos. Utiliza el paquete fancyhdr
% para darle un estilo moderno pero serio.
%
% La plantilla se encuentra adaptada al español.
%
%%%%%%%%%%%%%%%%%%%%%%%%%%%%%%%%%%%%%%%%%%%%%%%%%%%%%%%%%%%%%%%%%%%%%%%%%

%%%%%%%%%%%%%%%%%%%%%%%%%%%%%%%%%%%%%%%%%%%%%%%%%%%%%%%%%%%%%%%%%%%%%%%%%
% Plantilla de Trabajo
% Modificación de una plantilla de Latex de Frits Wenneker para adaptarla 
% al castellano y a las necesidades de escribir informática y matemáticas.
%
% Editada por: Mario Román
%
% License:
% CC BY-NC-SA 3.0 (http://creativecommons.org/licenses/by-nc-sa/3.0/)
%%%%%%%%%%%%%%%%%%%%%%%%%%%%%%%%%%%%%%%%%%%%%%%%%%%%%%%%%%%%%%%%%%%%%%%%%

%%%%%%%%%%%%%%%%%%%%%%%%%%%%%%%%%%%%%%%%%%%%%%%%%%%%%%%%%%%%%%%%%%%%%%%%%
% Short Sectioned Assignment
% LaTeX Template
% Version 1.0 (5/5/12)
%
% This template has been downloaded from:
% http://www.LaTeXTemplates.com
%
% Original author:
% Frits Wenneker (http://www.howtotex.com)
%
% License:
% CC BY-NC-SA 3.0 (http://creativecommons.org/licenses/by-nc-sa/3.0/)
%
%%%%%%%%%%%%%%%%%%%%%%%%%%%%%%%%%%%%%%%%%%%%%%%%%%%%%%%%%%%%%%%%%%%%%%%%%


% Tipo de documento y opciones.
\documentclass[11pt, a4paper]{article}


%---------------------------------------------------------------------------
%   PAQUETES
%---------------------------------------------------------------------------

% Idioma y codificación para Español.
\usepackage[utf8]{inputenc}
\usepackage[spanish, es-tabla, es-lcroman, es-noquoting]{babel}
\selectlanguage{spanish} 
%\usepackage[T1]{fontenc}

% Fuente utilizada.
\usepackage{courier}    % Fuente Courier.
\usepackage{microtype}  % Mejora la letra final de cara al lector.

% Diseño de página.
\usepackage{fancyhdr}   % Utilizado para hacer títulos propios.
\usepackage{lastpage}   % Referencia a la última página.
\usepackage{extramarks} % Marcas extras. Utilizado en pie de página y cabecera.
\usepackage[parfill]{parskip}    % Crea una nueva línea entre párrafos.
\usepackage{geometry}            % Geometría de las páginas.

% Símbolos y matemáticas.
\usepackage{amssymb, amsmath, amsthm, amsfonts, amscd}
\usepackage{upgreek}

% Otros.
\usepackage{enumitem}   % Listas mejoradas.
\usepackage[hidelinks]{hyperref}
\usepackage{graphicx}   % Gráficos.


%---------------------------------------------------------------------------
%   OPCIONES PERSONALIZADAS
%---------------------------------------------------------------------------

% Redefinir letra griega épsilon.
\let\epsilon\upvarepsilon

% Formato de texto.
\linespread{1.1}            % Espaciado entre líneas.
\setlength\parindent{0pt}   % No indentar el texto por defecto.
\setlist{leftmargin=.5in}   % Indentación para las listas.

% Estilo de página.
\pagestyle{fancy}
\fancyhf{}
\geometry{left=3cm,right=3cm,top=3cm,bottom=3cm,headheight=1cm,headsep=0.5cm}   % Márgenes y cabecera.

% Ruta donde buscar gráficos
\graphicspath{{../Recursos/Plantillas/}, {Recursos/Plantillas/}}

% Redefinir entorno de demostración (reducir espacio superior)
%\makeatletter
%\renewenvironment{proof}[1][\proofname] {\vspace{-15pt}\par\pushQED{\qed}\normalfont\topsep6\p@\@plus6\p@\relax\trivlist\item[\hskip\labelsep\it#1\@addpunct{.}]\ignorespaces}{\popQED\endtrivlist\@endpefalse}
%\makeatother

% Aumentar el tamaño del interlineado
\linespread{1.3}


%---------------------------------------------------------------------------
%   COMANDOS PERSONALIZADOS
%---------------------------------------------------------------------------

% Valor absoluto: \abs{}
\providecommand{\abs}[1]{\lvert#1\rvert}    

% Fracción grande: \ddfrac{}{}
\newcommand\ddfrac[2]{\frac{\displaystyle #1}{\displaystyle #2}}

% Texto en negrita en modo matemática: \bm{}
\newcommand{\bm}[1]{\boldsymbol{#1}}

% Línea horizontal.
\newcommand{\horrule}[1]{\rule{\linewidth}{#1}}

% Letras de conjuntos
\newcommand{\R}{\mathbb{R}}
\newcommand{\N}{\mathbb{N}}

% Sucesiones
\newcommand{\xn}{\{x_n\}}
\newcommand{\fn}{\{f_n\}}


%---------------------------------------------------------------------------
%   CABECERA Y PIE DE PÁGINA
%---------------------------------------------------------------------------

% Cabecera del documento.
\renewcommand\headrule{
	\begin{minipage}{1\textwidth}
		\hrule width \hsize 
	\end{minipage}
}

% Texto de la cabecera.
\lhead{\subject}  % Izquierda.
\chead{}            % Centro.
\rhead{\docauthor}    % Derecha.

% Pie de página del documento.
\renewcommand\footrule{                                 
	\begin{minipage}{1\textwidth}
		\hrule width \hsize   
	\end{minipage}\par
}

% Texto del pie de página.
\lfoot{}                                                 % Izquierda
\cfoot{}                                                 % Centro.
\rfoot{Página\ \thepage\ de\ \protect\pageref{LastPage}} % Derecha.


%---------------------------------------------------------------------------
%   ENTORNOS PARA MATEMÁTICAS
%---------------------------------------------------------------------------

% Nuevo estilo para definiciones.
\newtheoremstyle{definition-style} % Nombre del estilo.
{10pt}               % Espacio por encima.
{10pt}               % Espacio por debajo.
{}                   % Fuente del cuerpo.
{}                   % Identación.
{\bf}                % Fuente para la cabecera.
{.}                  % Puntuación tras la cabecera.
{.5em}               % Espacio tras la cabecera.
{\thmname{#1}\thmnumber{ #2}\thmnote{ (#3)}}     % Especificación de la cabecera (actual: nombre en negrita).

% Nuevo estilo para notas.
\newtheoremstyle{remark-style} 
{10pt}                
{10pt}                
{}                   
{}                   
{\itshape}          
{.}                  
{.5em}               
{}                  

% Nuevo estilo para teoremas y proposiciones.
\newtheoremstyle{theorem-style}
{10pt}                
{10pt}                
{\itshape}           
{}                  
{\bf}             
{.}                
{.5em}               
{\thmname{#1}\thmnumber{ #2}\thmnote{ (#3)}}                   

% Nuevo estilo para ejemplos.
\newtheoremstyle{example-style}
{10pt}                
{10pt}                
{}                  
{}                   
{\scshape}              
{:}                 
{.5em}               
{}                   

% Teoremas, proposiciones y corolarios.
\theoremstyle{theorem-style}
\newtheorem{nth}{Teorema}[section]
\newtheorem*{nprop}{Proposición}
\newtheorem{ncor}{Corolario}[section]

% Definiciones.
\theoremstyle{definition-style}
\newtheorem*{ndef}{Definición}

% Notas.
\theoremstyle{remark-style}
\newtheorem*{nota}{Nota}

% Ejemplos.
\theoremstyle{example-style}
\newtheorem{ejemplo}{Ejemplo}[section]

% Listas ordenadas con números romanos (i), (ii), etc.
\newenvironment{nlist}
{\begin{enumerate}
\renewcommand\labelenumi{(\emph{\roman{enumi})}}}
{\end{enumerate}}

% División por casos con llave a la derecha.
\newenvironment{rcases}
  {\left.\begin{aligned}}
  {\end{aligned}\right\rbrace}


%---------------------------------------------------------------------------
%   PÁGINA DE TÍTULO
%---------------------------------------------------------------------------

% Título del documento.
\newcommand{\subject}{Análisis Matemático II}

% Autor del documento.
\newcommand{\docauthor}{Doble Grado de Informática y Matemáticas}

% Título
\title{
  \normalfont \normalsize 
  \textsc{Universidad de Granada} \\ [25pt]    % Texto por encima.
  \horrule{0.5pt} \\[0.4cm] % Línea horizontal fina.
  \huge \subject\\ % Título.
  \horrule{2pt} \\[0.5cm] % Línea horizontal gruesa.
}

% Autor.
\author{\Large{\docauthor}}

% Fecha.
\date{\vspace{-1.5em} \normalsize Curso 2016/17}


%---------------------------------------------------------------------------
%   COMIENZO DEL DOCUMENTO
%---------------------------------------------------------------------------

\begin{document}

\maketitle  % Título.
\vfill
\begin{center}
\includegraphics{by-nc-sa}  % Licencia.
\end{center}
\newpage
\tableofcontents    % Índice
\newpage


%--------------------------------------------------------------------------------
%   Introducción.
%--------------------------------------------------------------------------------

\section*{Introducción.}

%%% FALTA COMPLETAR LA INTRODUCCIÓN ---------------------------------------------

\newpage

%--------------------------------------------------------------------------------
%   Convergencia uniforme y puntual.
%--------------------------------------------------------------------------------

\section{Convergencia uniforme y puntual}
De igual manera que tratamos con sucesiones de puntos de $\R^N$, podemos hacerlo con sucesiones de funciones. Dado $\emptyset \ne A \subseteq \R^N$, podemos tomar para cada $n \in \mathbb{N}$ una función $f_n : A \to \R^M$, y formar así una sucesión de funciones que notaremos $\{f_n\}$. El conjunto de funciones de $A$ en $\mathbb{R}^M$ lo denotaremos por $\mathcal{F}(A,\mathbb{R}^M)$. 

\begin{ndef}[Convergencia punto a punto] Diremos que una sucesión de funciones $\fn$ converge puntualmente a una función $f\in \mathcal{F}(A,\mathbb{R}^M)$ si $\forall x \in A \ \ \{f_n(x)\} \rightarrow f(x)$. Esto es, si se verifica lo siguiente:
	\[
		\forall x\in A\ \forall \epsilon > 0\ \exists n_0 \in \mathbb{N}: n \ge n_0 \Rightarrow |f_n(x)-f(x)| < \epsilon 
	\]
	
En ocasiones denotaremos la convergencia puntual como $\fn \xrightarrow {c.p} f$. 
\end{ndef}


\begin{ndef}[Convergencia uniforme] Diremos que $\fn$ converge uniformemente a una función $f \in \mathcal{F}(A,\mathbb{R}^M)$ si se verifica:
	\[
	\forall \epsilon>0 \quad \exists n_0 \in \mathbb{N}\ : n\ge n_0\implies \abs{f_n(x)-f(x)} < \epsilon \quad \forall x \in A
	\]
\end{ndef}

En ocasiones denotaremos la convergencia uniforme como $\fn \xrightarrow {c.u} f$.

\begin{nota}
Aunque ambas definiciones son muy parecidas, hay una diferencia clave. En la convergencia puntual, el valor de $n_0$ puede depender tanto de $\epsilon$ como de $x$. Sin embargo, en la convergencia uniforme, exigimos que $n_0$ sea válido para cualquier $x$.
\end{nota}
 
\begin{nprop}
	Si $\fn\to f \text{ uniformemente} \implies \fn \to f \text{ puntualmente}$. 
	\begin{nota}
	El recíproco no es cierto en general. Sin embargo, ambos conceptos son equivalentes si el conjunto de definición de las $f_n$ es un conjunto finito. En efecto, si $A$ es finito, y $\{f_n(x)\} \xrightarrow {c.p} f(x) \ \forall x \in A$, entonces tomamos $n_0$ como el máximo de los $n_0(x)$ que nos da la convergencia puntual en cada $x \in A$, y $\fn$ converge uniformemente en $A$.  
\end{nota}
\end{nprop}

La necesidad del concepto de convergencia uniforme se aprecia bien en el siguiente teorema, junto con los ejemplos que aparecen a continuación. Normalmente, las funciones con las que trabajemos serán continuas.

\begin{nth}
	Sean $\emptyset \ne A \subseteq \R^N$, $f_n \in \mathcal{C}(A,\R^M)\ \forall n\in \N$. Entonces, se tiene que
	\[
	\fn \to f \text{ uniformemente } \implies f \text{ es continua}
	\]
\end{nth}

\begin{proof}
	Fijamos $a\in A$. Dado $\epsilon>0$, como $\fn\to f$ uniformemente,
	\[
	\exists K>0:\ n\ge K \implies \abs{f_n(y)-f(y)}<\dfrac{\epsilon}{3}\ \forall y\in A \implies \begin{cases}
	\abs{f_n(a)-f(a)}<\dfrac{\epsilon}{3}\\
	\abs{f_n(x)-f(x)}<\dfrac{\epsilon}{3}
\end{cases}
	\]
	
	Como además, $f_n$ es continua para cada $n\in \N$,
	\[
	\exists \delta>0:\ \begin{rcases}
	\abs{x-a}<\delta\\
	x\in A
\end{rcases} \implies \abs{f_n(x)-f_n(a)} < \dfrac{\epsilon}{3}
	\] 
	
	Entonces,
	\[
	\exists \delta>0:\ \begin{rcases}
	\abs{x-a}<\delta\\
	x\in A
\end{rcases} \implies \abs{f(x)-f(a)} \le \abs{f(x)-f_n(x)} + \abs{f_n(x)-f_n(a)}+\abs{f_n(a)-f(a)} < \epsilon
	\]
	
	Hemos probado que dado $\epsilon > 0$, $\exists \delta>0:\ \begin{rcases}
	\abs{x-a}<\delta\\
	x\in A
	\end{rcases} \implies \abs{f(x)-f(a)}<\epsilon$, por tanto, f es continua. 
\end{proof}

Veamos algunoss ejemplos de sucesiones de funciones:

\begin{ejemplo}
	\[
		f_n : [0,1] \to \R,\ f_n(x) = \begin{cases}
	         0 & \text{si } x\ge \dfrac{1}{n}\\
	         -nx+1 & \text{si } 0\le x < \dfrac{1}{n}
\end{cases}
	\]
\end{ejemplo}

\begin{ejemplo}
	\[
	f_n : [0,1] \to \R,\ f_n(x) = x^n
	\]
\end{ejemplo}

\begin{ejemplo}
	\[
	f_n : [0,1] \to \R,\ f_n(x) = \dfrac{sen(nx)}{n}
	\]
\end{ejemplo}


Vamos a estudiar la convergencia puntual de la sucesión del ejemplo (1.1):

Primero, fijamos $x\in (0,1]$. Entonces, existe $n\in \N$ tal que $x \ge \dfrac{1}{n}$, luego $f_n(x) = 0$. Por otra parte, $f_n(0) = 1\ \forall n\in \N$. Concluimos que

\[
	\fn\to f = \begin{cases}
	1 & \text{ si } x=0\\
	0 & \text{ en otro caso}
\end{cases}
\]

Observamos que la convergencia puntual no preserva la continuidad de las funciones. Esto implicaría que, con esta definición de convergencia, el espacio de funciones continuas en un conjunto no sería cerrado. Además, podemos comprobar que $\fn$ no converge uniformemente a $f$, pues en caso de hacerlo $f$ debería ser continua, por el \textit{Teorema 1.1}.

Ahora estudiemos la convergencia uniforme del ejemplo (1.3):

\[
\forall\epsilon>0\ \exists K>\dfrac{1}{\epsilon}:\ \ n\ge K \implies \dfrac{\abs{sen(nx)}}{n} \le \dfrac{1}{n} \le \dfrac{1}{K} < \epsilon
\]

Vemos que converge uniformemente a cero. Lo importante para esta demostración, y lo que lo será en la mayoría de los casos de convergencia uniforme, es que podemos encontrar un $\epsilon_n$ (en este caso $\frac{1}{n}$), que no dependa de $x$, tal que $\abs{f_n(x)-f(x)} < \epsilon_n$.

\begin{nprop}[Criterio de Cauchy]
	Sea $A \subseteq \mathbb{R}^N$, y sean $f_n: A \longrightarrow \mathbb{R}^M \ \forall n \in \mathbb{N}$. Entonces: $$\fn \xrightarrow {c.u} f \iff \forall \epsilon > 0\ \exists n_0 \in \mathbb{N}:\ m,n \ge n_0 \Rightarrow |f_n(x) - f_m(x)| < \epsilon\ \ \forall x \in A$$
	
\begin{proof} \hfill \\
\boxed{\Rightarrow}	 $\fn \xrightarrow {c.u} f \implies \forall \epsilon > 0\ \exists n_0 \in \mathbb{N}:\ n \ge n_0 \Rightarrow |f_n(x) - f(x)| < \frac{\epsilon}{2}\ \forall x \in A$. Entonces, dados $m,n \ge n_0$ se tiene que: $$ |f_n(x) - f_m(x)| \le |f_n(x) - f(x) + |f_m(x) - f(x)| < \frac{\epsilon}{2} + \frac{\epsilon}{2} = \epsilon$$

\boxed{\Leftarrow} Sea $x \in A$ fijo. Entonces, es claro que $\{f_n(x)\} \subseteq \mathbb{R}^M$ es una sucesión de Cauchy. Como $\mathbb{R}^M$ es completo, tenemos que $\{f_n(x)\} \xrightarrow {c.p} f(x)\ \forall x \in A$. Ahora, tomando límite cuando $m \to \infty$ en la expresión de la hipótesis, y teniendo en cuenta que el último $<$ se transforma en $\le$: $$\forall \epsilon > 0\ \exists n_0 \in \mathbb{N}:\ n \ge n_0 \implies |f_n(x) - f(x)| \le \frac{\epsilon}{2} < \epsilon\ \ \forall x \in A$$

Es decir, $\fn \xrightarrow {c.u} f$.
\end{proof}
\end{nprop}

\subsection{El espacio de funciones continuas}

Ya sabemos que dado $A\subseteq \mathbb{R}^N$ compacto, el espacio $(\mathcal{C}(A,\mathbb{R}^M), ||\cdot||_{\infty})$ es un espacio normado, donde la norma del máximo o \textit{norma uniforme} se define así: $$||f||_{\infty} := m\acute{a}x \{ |f(x)|: x \in A\} = \underset{x\in A}{m\acute{a}x} \ |f(x)|$$

\begin{nprop} En el espacio $\mathcal{C}(A,\mathbb{R}^M)$, con $A$ compacto, la convergencia de sucesiones equivale a la convergencia uniforme, esto es: $$\fn \rightarrow f\ en\ \mathcal{C}(A,\mathbb{R}^M) \iff \fn \xrightarrow {c.u} f$$

\begin{proof} \hfill \\
$||f_n - f ||_{\infty} \rightarrow 0 \iff \underset{x \in A}{m\acute{a}x} \ |f_n(x) - f(x)| \rightarrow 0 \iff \forall \epsilon > 0\ \exists n_0 \in \mathbb{N}: n \ge n_0 \Rightarrow \underset{x \in A}{m\acute{a}x} \ |f_n(x) - f(x)| < \epsilon \iff |f_n(x) - f(x)| < \epsilon\  \forall x \in A \iff \fn \xrightarrow {c.u} f.$
\end{proof}
\end{nprop}

\begin{nth}[$\mathcal{C}(A,\R^M)$ es completo] En el espacio $\mathcal{C}(A,\mathbb{R}^M)$, con $A$ compacto, también ser sucesión de Cauchy equivale a la convergencia uniforme, esto es: $$\fn\ es\ de\ Cauchy\ en\ \mathcal{C}(A,\mathbb{R}^M) \iff \fn \xrightarrow {c.u} f $$

\begin{proof} El razonamiento es análogo al anterior, utilizando esta vez el criterio de Cauchy visto anteriormente.
\end{proof}
\end{nth}

Es importante recalcar que estamos suponiendo que el subconjunto $A$ es compacto. Podemos extender el resultado anterior, considerando el espacio $(\mathcal{C}_B(A,\mathbb{R}^M), \|\cdot\|_{\infty})$, donde: 

$\mathcal{C}_B(A,\mathbb{R}^M) := \{ f:A \longrightarrow \mathbb{R}^M: \text{$f$ es continua y acotada}\}$

$\|f\|_{\infty} := \underset{x \in A}{sup} \ \ |f(x)|$

\begin{nprop} El espacio $\mathcal{C}_B(A, \mathbb{R}^M)$ es un espacio de Banach, es decir, es un espacio normado y completo.

\begin{proof}
	Puede consultarse en \textit{[Ejercicios.pdf, pág 2]}.
\end{proof}

\end{nprop}


\begin{nota}
Si $A$ es compacto, entonces  $\mathcal{C}_B(A,\mathbb{R}^M) = \mathcal{C}(A, \mathbb{R}^M)$.
\end{nota}

Veamos ahora dos teoremas que relacionan el concepto de convergencia uniforme con los conceptos de derivación e integración.

\begin{nth} Sean $f, f_n \in \mathcal{C}([a,b],\mathbb{R})$ tales que $\fn \xrightarrow {c.u} f$. Entonces, $$\left\{\int_a^b f_n\right\} \xrightarrow {c.u} \int_a^b f $$

Equivalentemente, podemos intercambiar la integral con el límite: $$\lim_{n\rightarrow \infty} \int_a^b f_n = \int_a^b \lim_{n\rightarrow \infty} f_n$$

\begin{proof} Dado $\displaystyle \epsilon >0\ \exists n_0 \in \mathbb{N}: n \ge n_0 \Rightarrow |f_n(x)-f(x)| < \frac{\epsilon}{b-a}\ \forall x \in [a,b]$. Entonces, $$\left| \int_a^b f_n - \int _a^b f \right| = \left| \int_a^b (f_n - f) \right| \le \int_a^b |f_n - f| < \int_a^b \frac{\epsilon}{b-a} = \epsilon$$
\end{proof}
\end{nth}


\begin{nth}
	Sea $f_n\in \mathcal{C}^1((a,b)) \ \forall n \in \mathbb{N}$, y sean $f,g \in \mathcal{C}((a,b)$. Supongamos que $ \{f_n(x)\} \to f(x)$ c.p. $\forall x \in (a,b)$, y supongamos también que $\{f_n'\} \xrightarrow{c.u.} g $ en (a,b). Entonces, $f \in \mathcal{C}^1((a,b))$, y $f' = g$.
	
\begin{proof}
	Elegimos primero un $x_0 \in (a,b)$ fijo. Entonces, $\{f_n(x_0)\} \to f(x_0)$ por hipótesis.
	Como $f_n'$ es continua, entonces, por el teorema fundamental del cálculo:
	\[
	f_n(x) = f_n(x_0) + \int_{x_0}^x f_n'(t)dt
	\]
	Ahora, como  $\{f_n'\} \xrightarrow{c.u.} g $ en el intervalo cerrado de extremos $x_0$ y $x$, por el \textit{Teorema 1.3} tenemos que $\{f_n(x)\} \to G(x)\ \forall x \in (a,b)$, donde:
	\[
	    G(x) = f(x_0) + \int_{x_0}^x g(t)dt
	\]
	Es decir, $\{f_n(x)\}$ converge puntualmente en $(a,b)$ a $G(x)$. Ahora, $G$ es de clase 1 por ser $g(t)$ continua, y además, tenemos que $G(x_0) = f(x_0)$.
	Por otro lado, es claro que $G' = g$.
	
	Pero también $\{f_n(x)\} \xrightarrow {c.p} f(x)$ por hipótesis, por lo que necesariamente $\forall x \in (a,b)$ \\ $G(x) = f(x)$, esto es, $f \in \mathcal{C}^1((a,b))$ y $f' = G' = g$.
\end{proof}
\end{nth}

\begin{ncor}
	Sea $f_n\in \mathcal{C}^1((a,b)) \ \forall n \in \mathbb{N}$, y sean $f,g \in \mathcal{C}((a,b)$. Supongamos que $\exists \alpha \in (a,b)$ tal que $ \{f_n(\alpha)\}$ es convergente, y supongamos también que $\{f_n'\} \xrightarrow{c.u.} g $ en (a,b). Entonces, $\exists f: (a,b) \to \mathbb{R}$ tal que $\fn \to f$ uniformemente. Además, $f \in \mathcal{C}^1(a,b)$, y  $f'(x) = g(x)\ \forall x \in (a,b)$.

\begin{proof}
	Probaremos únicamente que $\fn \to f$ uniformemente. El resto de la tesis se sigue de aplicar el \textit{Teorema 1.4}.
	
Sea $\epsilon > 0$. Por un lado, como $\{f'_n\}$ converge uniformemente, aplicamos el \textit{Criterio de Cauchy} para obtener un $k_1 \in \mathbb{N}$ tal que: $$p,q \ge k_1 \implies |f'_p(y) - f'_q(y)| < \frac{\epsilon}{2 (b-a)}$$
Del mismo modo, como $\{f_n(\alpha)\}$ es una sucesión de vectores convergente, es de Cauchy, y obtenmos un $k_2 \in \mathbb{N}$ tal que: $$p,q \ge k_2 \implies |f_p(\alpha) - f_q(\alpha)| < \frac{\epsilon}{2}$$
Ahora, defino $k := k_1 + k_2$, y teniendo en cuenta el \textit{Teorema del Valor Medio}, se tiene que, dado $x \in (a,b)$: $$|f_p(x) - f_q(x)| \le |f_p(x) - f_p(\alpha) + f_p(\alpha) - f_q(\alpha) + f_q(\alpha) - f_q(x)|$$ $$\le |f_p(x) - f_p(\alpha)| + |f_q(\alpha) - f_q(x)| + |f_p(\alpha) - f_q(\alpha)| \le \left|x-a \right|\left| f'_p(z_x) - f'_q(z_x) \right| + |f_p(\alpha) - f_q(\alpha)|$$ $$< (b-a) \cdot \frac{\epsilon}{2(b-a)} + \frac{\epsilon}{2} = \epsilon$$

donde $z_x$ está en el intervalo abierto de extremos $a$ y $x$. Observemos que en este razonamiento, el valor de $m$ \textbf{no depende de $x$}, sino exclusivamente de quién sea $\epsilon$. Por tanto, por el \textit{Criterio de Cauchy}, se tiene que $\{f_n\}$ converge uniformemente en $(a,b)$.
\end{proof}
	
\end{ncor}


\subsection{Conjuntos compactos de $\mathcal{C}(A, \mathbb{R}^M$)}

En este apartado, nuestro objetivo será caracterizar los conjuntos compactos de $\mathcal{C}(A,\mathbb{R}^M)$. Para ello, necesitamos primero definir ciertos conceptos el espacio de funciones continuas. Salvo que se especifique lo contrario, $A$ será un subconjunto compacto de $\mathbb{R}^N$.

\begin{ndef}[Conjunto acotado]
	Sea $A \in \mathbb{R}^N$ compacto, $B \subseteq \mathcal{C}(A,\mathbb{R}^M)$. Se dice que $B$ es acotado si $\exists M > 0: \ \|f\|_{\infty} \le M \ \ \forall f \in B$.
\end{ndef}

\begin{ndef}[Conjunto equicontinuo]
    Sea $A \in \mathbb{R}^N$ compacto, $B \subseteq \mathcal{C}(A,\mathbb{R}^M)$. Se dice que $B$ es equicontinuo si $$\forall \epsilon > 0\ \exists \delta > 0:\ \begin{rcases} |x-y| < \delta\\ x,y \in A \end{rcases} \implies |f(x) - f(y)| < \epsilon \quad \forall f \in B$$
\end{ndef}

\begin{nota}
	Este concepto es muy parecido al de función uniformemente continua. De hecho, como $A$ es compacto y todas las $f$ son continuas, sabemos por el \textit{Teorema de Heine-Cantor} que todas ellas son uniformemente continuas. Sin embargo, para poder afirmar que $B$ es equicontinuo, es necesario que el número $\delta$ sea válido sea quien sea $f$. 
\end{nota}

\begin{nprop} Sea $B \subseteq \mathcal{C}(A,\mathbb{R}^M)$. Si B es finito, entonces es equicontinuo.

\begin{proof}
	Sabemos que $\forall f \in B,\ f$ es uniformemente continua. Como $B$ es finito, tenemos un conjunto finito $\{\delta_i > 0 : f_i \ \text{es uniformemente continua con} \ \delta_i\}$. Concluimos tomando $\delta := \underset{i}{\text{mín}} \ \delta_i$.
\end{proof}
\end{nprop}

El interés en ver cuáles son los conjuntos compactos de $\mathcal{C}(A,\mathbb{R}^M)$ proviene, entre otras cosas, del hecho de que hay conjuntos cerrados y acotados que no son compactos. El siguiente ejemplo lo pone de manifiesto.

\begin{ejemplo} En $\mathcal{C}(A,\mathbb{R}^M)$, ninguna bola cerrada y centrada en $0$ es compacta.

%%% NIGNUNA BOLA, O SOLO LAS CENTRADAS EN EL ORIGEN? ---------------------------

Para verlo, sea $r>0$, y consideremos la sucesión de funciones $\fn$ definidas como $f_n(x) = rx^n \ \ \forall x \in [0,1]$. Ya sabemos que la sucesión converge puntualmente a la función $f: \mathbb{R} \to \mathbb{R}$ dada por $$f(x) = \begin{cases} 0 & si\ 0 \le x < 1\\ r & si\ x = 1 \end{cases}$$

Pero $f$ no es continua, por lo que $\fn$ no converge uniformemente en $[0,1]$. De hecho, $\fn$ no puede tener ninguna sucesión parcial que converja uniformemente, pues en otro caso, el límite uniforme de la sucesión parcial sería $f$, que no es continua.

Por otro lado, $\|f_n\|_{\infty} = \text{máx} \ \{ |f_n(x)|:\ x \in [0,1]\} = r$, y entonces $f_n \in \overline{B_{\infty}(0,r)}$. 

Recordemos que en este espacio, ser convergente era equivalente a la convergencia uniforme. Por tanto, combinando ambas informaciones, hemos encontrado una sucesión de funciones $\fn \subseteq \overline{B_{\infty}(0,r)}$ tal que no es posible extraer una subsucesión convergente a un punto de $\overline{B_{\infty}(0,r)}$. Esto prueba que $\overline{B_{\infty}(0,r)}$ no es compacta.
\end{ejemplo}

\begin{ndef}[Sucesión equicontinua]
	Se dice que una sucesión $\fn \subseteq \mathcal{C}(A,\mathbb{R}^M)$ es equicontinua si el conjunto $\{ f_n: n \in \mathbb{N}\}$ es equicontinuo. 
\end{ndef}

Si recordamos la prueba en $\mathbb{R}^N$ de que los conjuntos compactos son los cerrados y acotados, el \textit{Teorema de Bolzano-Weierstrass} era una herramienta clave. Veamos ahora un resultado equivalente en $\mathcal{C}(A,\mathbb{R}^M)$.

\begin{nth}[Arzelà-Ascoli] Toda sucesión de funciones en $\mathcal{C}(A,\mathbb{R}^M)$ que sea equicontinua y acotada admite una sucesión parcial convergente.

%%% COMPLETAR DEMOSTRACIÓN (DIAPOSITIVAS) --------------------------------------

\begin{proof} Dividiremos la demostración en dos pasos.

\underline{Paso 1.} Vamos a probar que $\exists \sigma : \mathbb{N} \to \mathbb{N}$ estrictamente creciente tal que $\{f_\sigma(n)\} \to f$ puntualmente en $\mathbb{Q}\cap A$.  \textbf{COMPLETAR.}

\underline{Paso 2.}
\end{proof}
	
\end{nth}

Como conclusión de este capítulo, presentamos el siguiente corolario, que nos sirve para caracterizar los conjuntos compactos.

\begin{ncor}[Caracterización de compactos] Sea $B \in \mathcal{C}(A,\mathbb{R}^M)$. Entonces: $$B\ \text{es compacto} \iff B \ \text{es cerrado, acotado y equicontinuo}$$

\begin{proof} \hfill \\
	\boxed{\Rightarrow} Ya sabemos que un conjunto compacto de un espacio métrico es cerrado y acotado. Veamos que además es equicontinuo.
	
%%% COMPLETAR DEMOSTRACIÓN (EJERCICIO PRADO).

    \boxed{\Leftarrow} Sea $\fn \subseteq B$. Es claro que $\fn$ es acotada y equicontinua, y aplicando el \textit{Teorema de Arzelá-Ascoli}, tenemos que $\exists \{f_{\sigma(n)}\} \to f \in \overline{B}$. Pero $B$ es cerrado, por lo que $\overline{B} = B$. Obtenemos así que de cualquier sucesión de puntos de $B$ es posible obtener una sucesión parcial convergente a un punto de $B$, por lo que $B$ es compacto.
\end{proof}
	
\end{ncor}

%--------------------------------------------------------------------------------
%   Series de funciones.
%--------------------------------------------------------------------------------
\newpage

\section{Series de funciones}

\begin{ndef}[Sumas parciales]
Sea $\Omega\subset \mathbb R^N$ un conjunto. 
Dada una sucesi\'on $\{f_n \}$ de funciones $f_n : \Omega \longrightarrow \mathbb R^M$, se llama sumas parciales $S_k$ a la funci\'on definida en $\Omega$ mediante
 \[ 
 S_k (x) = \sum_{n=1}^{k} f_n(x),\ \ x\in\Omega.
 \] 
\end{ndef}

\begin{ndef}[Serie de funciones]
Llamamos serie  de funciones $\sum_{n\geq 1} f_n$ al par formado por la sucesiones  $\{ f_n \}$ y $\{S_n\}$. Llamaremos a $f_n$ el t\'ermino general de la serie de funciones $\displaystyle\sum_{n\geq 1} f_n$.
\end{ndef}

\begin{ndef}[Convergencia de series de funciones]
El car\'acter de la convergencia de la serie $\sum_{n\geq 1} f_n$ de funciones se basa en la correspondiente convergencia de la sucesi\'on de funciones $\{S_n\}$ formada por las sumas parciales $S_n$. Por tanto, tenemos de nuevo dos tipos de convergencia.

Si $S:\Omega\longrightarrow \mathbb R^M$ es una funci\'on y $A\subset \Omega$, diremos 
$$\sum_{n\geq 1} f_n  \mbox{ es convergente puntualmente en } A \Leftrightarrow \{S_n \} \longrightarrow S\ c.p.\ en\ A.$$

$$\sum_{n\geq 1} f_n  \mbox{ es convergente uniformente en }A \Leftrightarrow \{S_n(x)\} \longrightarrow S(x) \mbox{ c.u. en }A.$$ 

Diremos que $S$ es la suma de la serie de funciones: 
\[
S=\displaystyle \sum_{n=1}^{\infty} f_n \mbox{ puntualmente (resp. uniformemente) en $A$.}
\]
\end{ndef}

\begin{nprop}
	Sea $\sum_{n \ge 0} f_n$ una serie de funciones definida en un conjunto $A \subseteq \mathbb{R}^N$. Si dicha serie converge uniformemente en un subconjunto $B \subseteq A$, entonces la sucesión $\fn$ converge a 0 uniformemente en $B$.
	
\begin{proof}
Es claro que $\forall n \in \N,\ f_n = S_{n+1} - S_n$. Aplicando límites en la igualdad anterior, y sabiendo que $\sum_{n\ge 0} f_n$ converge uniformemente en $B$, tenemos que: 

$$\lim_{n \to \infty} \{f_n(x)\} = \lim_{n \to \infty} \{S_n(x)\} - \lim_{n \to \infty} \{S_{n + 1}(x)\} = S(x) - S(x) = 0\quad \forall x \in B$$
\end{proof}
\end{nprop}

Puesto que  la convergencia de la serie $\sum_{n\geq 1} f_n$ de funciones se basa en la correspondiente convergencia de la sucesi\'on de funciones $\{S_n\}$ formada por las sumas parciales, tenemos una relación directa de la convergencia de series con la continuidad, derivación e integración.


\begin{nth}
Si una serie  $\sum_{n\geq 1} f_n $ de funciones continuas $f_n$ es uniformemente convergente, entonces la funci\'on suma $\displaystyle \sum_{n=1}^{\infty} f_n $ es una funci\'on continua.

\begin{proof}
Como la serie es uniformemente convergente, tenemos que $\{S_n\} \to S$ uniformemente. Como las $f_n$ eran continuas, las funciones $S_n$ son continuas, por ser suma finita de funciones continuas. Por tanto, aplicando el teorema \textit{Teorema 1.1} a la sucesión de funciones $\{S_n\}$, tenemos que $S = \sum_{n=1}^{\infty} f_n$ es continua.
\end{proof}
\end{nth}


\begin{nth}
Si la serie $\sum_{n\geq 1} f_n $ de funciones continuas $f_n:[a,b]\longrightarrow \mathbb R$ es uniformemente convergente, entonces 
$
\displaystyle \int_a^b \sum_{n=1}^{\infty} f_n (x) \, dx =  \sum_{n=1}^{\infty} \int_a^b f_n (x) \, dx  .
$

\begin{proof}
Siguiendo la idea del teorema anterior, aplicamos el \textit{Teorema 1.3} a la sucesión de funciones $\{S_n\}$, y obtenemos que:

$$\int_a^b \lim_{n \to \infty} S_n = \lim_{n \to \infty} \int_a^b S_n \Rightarrow \int_a^b \sum_{n=1}^{\infty} f_n = \lim_{k \to \infty} \sum_{n=1}^k \int_a^b f_n = \sum_{n =1}^{\infty} \int_a^b f_n$$
\end{proof}
\end{nth}


\begin{nth}
Si para $f_n\in \mathcal{C}^1(a,b)$  la serie  $\sum_{n\geq 1} f_n $ es  c.p.  en $(a,b)$ y  la serie  $ \sum_{n\geq 1} f'_n $ es  c.u. en (a,b), entonces la suma $\displaystyle \sum_{n=1}^{\infty} f_n $ es una funci\'on de  $\mathcal{C}^1(a,b)$ con 
$
\displaystyle \left( \sum_{n=1}^{\infty} f_n \right)' =  \sum_{n=1}^{\infty}  f'_n  .
$

\begin{proof}
Claramente $S_n \in \mathcal{C}^1(a,b)\ \forall n \in \mathbb{N}$, y además tenemos que $\{S_n\}\ c.p.$ en $(a,b)$, y $\{S'_n\}\ c.u.$ en $(a,b)$, pues la derivada de una suma finita es la suma de las derivadas. Bajo estas hipótesis, podemos aplicar a $\{S_n\}$ el \textit{Teorema 1.4}, obteniendo que $S = \sum_{n=1}^{\infty} f_n \in \mathcal{C}^1(a,b)$, y además:

$$S' = \left( \sum_{n=1}^{\infty} f_n \right)' = \lim_{n \to \infty} \{S'_n\} = \sum_{n=1}^{\infty} f'_n$$
\end{proof}
\end{nth}

\subsection{Criterios de convergencia para series de funciones}

\begin{nth}[Criterio de Weierstrass]
	Si existen constantes $ M_{n} $ positivas tales que $ |f_{n}(x)| \leq M_{n}\ \ \forall x \in A$ y la serie de números reales $ \sum_{n \ge 1}M_{n} $ es convergente, entonces la serie de funciones $ \sum_{n \ge 1}f_{n} $ es convergente uniformemente (y absolutamente) en A.
	
\begin{proof}
	Por el criterio de Cauchy para la convergencia de una serie de números positivos (las sumas parciales son de Cauchy), se tiene que:
	$$ \forall \epsilon > 0\ \exists n_{o} \in \mathbb{N}:\ |S_{n+k} - S_{n-1}| = M_{n} + ... + M_{n+k} < \epsilon\ \  \forall n \geq n_{o}\ \forall k \in \mathbb{N} $$
	
	Usando que $ |f_{n}(x)| \leq M_{n}\ \forall x \in A $ y la desigualdad triangular obtenemos:
	$$ |f_{n}(x) + ... + f_{n+k}(x)| \leq |f_{n}(x)| + ... + |f_{n+k}(x)| \leq M_{n} + ... + M_{n+k} < \epsilon\ \ \forall x \in A $$
	
	Y por tanto:
	$$\forall \epsilon > 0\ \exists n_{o} \in \mathbb{N}:\ |f_{n}(x) + .... + f_{n+k}(x)| < \epsilon\ \  \forall n \geq n_{o}\ \forall k \in \mathbb{N}\ \forall x \in A $$
	
	Como la sucesión de sumas parciales de la serie de funciones $\sum_{n\ge 1} f_n$ converge uniformemente, queda probado que dicha serie converge uniformemente en A.
\end{proof}
\end{nth}

\begin{nth}[Criterio de Abel] Sean $ A \subseteq \mathbb{R}^n\  y \ \phi_n : A \rightarrow \mathbb{R}$ tales que $\{\phi_n\}$ es una sucesión decreciente de funciones, es decir, $\phi_{n+1}(x) \leq \phi_n (x)\ \forall x \in A.$ Supóngase que $\exists M$ tal que $|\phi_n (x)| \leq M\ \forall x \in A\ \forall n \in \mathbb{N}.$ Entonces, se tiene que
$$\sum_{n \ge 1} f_n(x)\ c. u.\ en\ A\, \implies \sum_{n \ge 1} \phi_n(x) f_n(x)\ c. u.\ en\ A.$$
\end{nth}

\begin{nth}[Criterio de Dirichlet] Sea $A \subseteq \mathbb{R}^n\  y \ \phi_n : A \rightarrow \mathbb{R}$ tales que $\{\phi_n\}$ es una sucesión decreciente de funciones, que converge a 0 uniformemente, con $\phi_n \ge 0 \ \forall n \in \mathbb{N}$. Sea $\sum_{n \ge 1} f_n$ una serie de funciones, y supóngase que $\exists M$ tal que las sumas parciales de dicha serie están uniformemente acotadas, esto es, $$\left| \sum_{k=1}^n f_k(x) \right| \le M \ \ \forall x \in A, \ \forall n \in \mathbb{N}$$ 

Entonces, la serie $\displaystyle \sum_{n\ge 1} \phi_n(x)f_n(x)$ converge uniformemente en $A$.
	
\end{nth}

Veamos un ejemplo de aplicación de estos criterios.

\begin{ejemplo}
	Probar que la serie $\displaystyle \sum_{n \ge 1} \frac{(-1)^n}{n} e^{-nx}$ converge uniformemente en $[0, \infty)$.
	
Sea $\phi_n(x) = e^{-nx}$. Para $x \ge 0$, $\phi_n$ es decreciente, y además $|\phi_n(x)| = | e^{-nx} |  = |\frac{1}{e^{nx}} |\le 1$. Por otro lado, sabemos que la serie alternada $ \sum_{n \ge 1} \frac{(-1)^n}{n}$ converge, por el \textit{Criterio de Leibnitz} para series de números reales. Por tanto, aplicando el \textit{Criterio de Abel}, la serie de partida converge uniformemente para $x \ge 0$.  
\end{ejemplo}

\subsection{Series de potencias}

Un caso interesante de series de funciones lo constituyen las series de potencias, esto es, aquellas series en las que las funciones son potenciales: $$\sum_{n \ge 0} a_n(x-a)^n, \quad a_n, a \in \mathbb{R} \ \ \forall n \in \mathbb{N}$$

Veamos una característica importante de las series de potencias, que nos ayuda a determinar su carácter de convergencia.

\begin{ndef}[Radio de convergencia] Dada una serie de potencias $\sum_{n \ge 0} a_n(x-a)^n$, definimos su radio de convergencia $R \in [0, \infty]$ como: $$
\frac{1}{R} =\underset{n\to\infty}{\mbox{lim\,sup\,}}  \sqrt[n]{|a_n|},$$
\end{ndef}

\begin{nota}
	En la definición anterior, permitimos que $R = \infty$, tomando como convenio que $\frac{1}{0} = \infty$ y $\frac{1}{\infty} = 0$.
\end{nota}

\begin{ndef}[Disco de convergencia] Dada una serie de potencias $\sum_{n \ge 0} a_n(x-a)^n$, definimos el disco de convergencia $D(a,R) := \{ x : |x-a| < R\}$. 
\end{ndef}

\begin{nota}
	Si $R = \infty$, entonces $D(a,R) = \mathbb{R}$.
\end{nota}

El siguiente teorema da sentido a la nomeclatura seguida en las dos definiciones anteriores.

\begin{nth}
La serie de potencias $\displaystyle\sum_{n\geq 0} a_n (x-a)^n$  
\begin{nlist}
\item converge uniformemente en $D(a,R')$ para cualquier $R'<R$ (en particular, converge absolutamente para todo $x\in D(a,R)$.
\item no converge si $x\not\in \overline{D(a,R)}$.
\end{nlist}

\begin{proof} \hfill \\
	\textit{(i)} Sea $R'<R$. Como $\underset{n\to\infty}{\mbox{lim\,sup\,}} \sqrt[n]{|a_n|} < 1/R'$, si elegimos $R''\in (R',R)$, deducimos que $\exists n_0\geq 0$ tal que $\sqrt[n]{|a_n|} \leq 1/R''$, $\forall n\geq n_0$.

Esto implica que
\[
\left| a_n (x-a)^n\right| = |a_n|\cdot|x-a|^n \leq \left(\frac{R'}{R''}\right)^n, \ \ \forall x\in D(a,R'), \ \forall n\geq n_0,
\]
y el primer apartado se concluye aplicando el criterio de Weierstrass.

\textit{(ii)} Basta observar que si la serie $\sum_{n\geq 0} a_n (x-a)^n$ es convergente en un punto $x\not= a$, entonces $\left\{ a_n (x-a)^n \right\} \to 0$, y tomando $\epsilon = 1$ en la definición de convergencia tenemos que:
\[
\exists n_0>0: \,  |a_n| |x-a|^n \leq 1, \ \forall n\geq n_0
\]
Y por tanto, $$\underset{n\to\infty}{\mbox{lim\,sup\,}} \sqrt[n]{a_n} \cdot |x-a| = |a_n| \cdot |x-a|^n =\frac{|x-a|}{R} \leq 1 \implies |x-a| \le R$$

Es decir, si la serie converge, necesariamente $x \in \overline{D(a,R)}.$
\end{proof}
\end{nth}

Por último, veamos un teorema sobre derivación e integración de series de potencias. Para ello, notemos que si $R$ es el radio de convergencia de la serie de potencias $\displaystyle\sum_{n\geq 0} a_n (x-a)^n$, podemos definir la función $S$ en $D(a,R)$ mediante: \[ 
S(x):= \sum_{n= 0}^\infty a_n (x-a)^n
, \ \ \forall x\in  D(a,R)
\] 

\begin{nth}
Sea $\displaystyle\sum_{n\geq 0} a_n (x-a)^n$ una serie de potencias. Entonces, se verifican las siguientes afirmaciones:	
\begin{nlist}
	\item La suma $S(x)=\displaystyle\sum_{n= 0}^\infty a_n (x-a)^n$ es una función $\mathcal{C}^\infty$ en el disco de convergencia $D(a,R)$.
\item $S'(x)= \displaystyle\sum_{n= 1}^\infty n a_n (x-a)^{n-1}$, esto es,
\[
\displaystyle\frac{d}{dx} \left( \sum_{n= 0}^\infty a_n (x-a)^n \right) =  \sum_{n= 1}^\infty \frac{d}{dx} \left( a_n (x-a)^n \right) .
\]
\item 
$\displaystyle \int_a^x S(y)\, dy =   \sum_{n= 0}^\infty \frac{a_n}{n+1} (x-a)^{n+1}$, esto es,
\[
\int_a^x  \sum_{n= 0}^\infty a_n (x-a)^n \, dx 
=  \sum_{n= 0}^\infty \int_a^x a_n (x-a)^n \, dx.
\]
\end{nlist}

%%% FALTA DEMOSTRAR EL TEOREMA ----------------------------------------------------
\end{nth}


\subsection{Funciones analíticas}

Nos preguntamos ahora si cualquier función $\mathcal{C}^{\infty}$ se puede escribir como suma de una serie de potencias, es decir, si dada $f \in \mathcal{C}^{\infty}(I)\ y\ a \in I, \ \exists \{a_n\}\ tal\ que\ f(x) = \sum_{n=0}^{\infty}a_n(x-a)^n$.


Haciendo uso del teorema de Taylor, podemos ver la siguiente condici\'on suficiente para que una funci\'on $f\in \mathcal{C}^\infty$ coincida con la suma de su serie de Taylor.

\begin{nth}
Si una funci\'on de clase $\mathcal{C}^\infty$ en un intervalo $I$ verifica 
\[
\exists M>0 \ : \  |f^{n)}(x)|\leq M, \ \ \forall x\in I
\]
entonces 
\[
f(x) = \sum_{n=0}^\infty \frac{f^{n)} (a)}{n!} (x-a)^n , \ \ \forall x\in I.
\]

\begin{proof}
	Fijamos $x \in I$, y observamos que las sumas parciales $S_k(x)$ de la serie son el polinomio de Taylor de orden $k$: 
\[
P_k(x) =  \sum_{n=0}^k \frac{f^{n)} (a)}{n!} (x-a)^n,
\]
Usando el teorema del resto, tendremos que $\exists c$ entre $x$ y $a$ tal que:
\[
|f(x)-S_k(x)| =\left| \frac{ f^{k+1)} (c)}{(k+1)!} (x-a)^{k+1}\right| \leq  M \frac{|x-a|^{k+1}}{(k+1)!} \to 0
\]

Es decir, la distancia entre $f$ y $S_n$ se hace tan pequeña como se quiera. Por tanto, hemos probado que $\displaystyle f(x) = \lim_{n \to \infty} \{S_n(x)\} = \sum_{n=0}^{\infty} \frac{f^{n)} (a)}{n!} (x-a)^n, \quad \forall x \in I$.
\end{proof}
\end{nth}
Como consecuencia de este teorema, podemos escribir algunas funciones elementales como la suma de su serie de Taylor, en el punto $a = 0$.

\begin{ejemplo}
	$e^x = \displaystyle \sum_{n=0}^\infty  \frac{x^n}{n!}$. 
\end{ejemplo}

\begin{ejemplo}
	$\sen x = \displaystyle \sum_{n=0}^\infty   \frac{(-1)^n x^{(2n+1)}}{(2n+1)!}$.
\end{ejemplo}
\begin{ejemplo}
	$\cos x = \displaystyle \sum_{n=0}^\infty   \frac{(-1)^n x^{2n}}{(2n)!}$.
\end{ejemplo}


Sin embargo, no todas las funciones $\mathcal{C}^{\infty}$ admiten una descomposición en sumas de funciones potenciales. Dada una serie de potencias, y usando el \textit{Teorema 2.7}, notemos lo siguiente:

$\displaystyle S'(a) = \sum_{n=1}^{\infty} na_n(a-a)^{n-1} = a_1 \quad ; \quad S''(a) = \sum_{n=2}^{\infty} n(n-1)a_n(a-a)^{n-2} = 2a_2$ \quad ; \quad \dots

En general, se demuestra por inducción que $S^{k)} (a) = k! \cdot a_k$. Es decir, fijado un $a$ y conocida la función $S$, se tiene que necesariamente $a_k = \frac{S^{k)}(a)}{k!}$.

Recordemos que si una función $f \in \mathcal{C}^{\infty}$ cumple que $f(x) = \sum_{n=0}^{\infty} a_n(x-a)^n\ \forall x \in A$, dicha función es justamente lo que nos hemos referido como la suma de la serie, $S$. Pero acabamos de ver que, en ese caso, los $a_n$ están perfectamente determinados, y resulta que $f$ se escribiría como su suma de Taylor: $$f(x) = \sum_{n=0}^{\infty} \frac{f^{n)}(a)}{n!} (x-a)^n$$

Por tanto, la pregunta que nos hacíamos equivale a preguntarse si toda función $\mathcal{C}^{\infty}$ se puede escribir como su suma de Taylor. La respuesta es que, en general, esto no es posible (ya vimos en el \textit{Teorema 2.8} que sí se cumple si todas las derivadas de $f$ están acotadas).

Un contraejemplo sería la función $f: \mathbb{R} \longrightarrow \mathbb{R}$ dada por: $$f(x) = \begin{cases}
	         e^{-1/x^2} & \text{si } x\ne 0\\
	         0 & \text{si } x=0
\end{cases}$$ Es claro que $f \in \mathcal{C}^{\infty}$ y que $f \ne 0$. Además, se tiene que $f(0) = 0$, y $f^{k)}(0) = 0\ \ \forall k \ge 1$. 

Si $f$ se escribiese como su suma de Taylor, tendríamos que, en el punto $a=0$: $$f(x) = \sum_{n=0}^{\infty} \frac{f^{n)}}{n!} \cdot x^n =  \sum_{n=0}^{\infty} \frac{0}{n!}\cdot x^n = 0\ \ \forall x \in \mathbb{R}$$

Lo cual es una contradicción, pues $f\ne 0$. Por tanto, $f$ no se puede escribir como suma de una serie de potencias.

Las funciones que \textit{sí} se pueden escribir como suma de una serie de potencias tienen propiedades muy interesantes. Tanto es así, que dichas funciones tienen un nombre propio.

\begin{ndef}[Función analítica] Decimos que una función $f \in \mathcal{C}^{\infty}(A)$ es analítica si se puede escribir como su suma de Taylor, es decir, si se cumple que: $$f(x) = \sum_{n=0}^{\infty} \frac{f^{n)}(a)}{n!}(x-a)^n\ \ \forall x \in A$$
\end{ndef}

\end{document}
